\documentclass[12pt, a4paper, oneside]{report}

% ─── Packages ───────────────────────────────────────────────
\usepackage[top=1.3in, bottom=1.3in, left=1in, right=1.3in]{geometry}
\usepackage{mathptmx}
\usepackage{setspace}
\usepackage{titlesec}
\usepackage{fancyhdr}
\usepackage{graphicx}
\usepackage{hyperref}
\usepackage{tocloft}
\usepackage{parskip}
\usepackage{enumitem}
\usepackage{booktabs}
\usepackage{array}
\usepackage{tabularx}
\usepackage{tikz}
\usetikzlibrary{calc}
\usepackage{eso-pic}
\usepackage{ragged2e}

\hypersetup{
    colorlinks=true,
    linkcolor=black,
    urlcolor=blue,
    citecolor=black
}

\onehalfspacing

% ─── Header/Footer ──────────────────────────────────────────
\pagestyle{fancy}
\fancyhf{}
\fancyfoot[C]{\small NMIET, Department of Computer Engineering 2025-26}
\fancyfoot[R]{\small \thepage}
\renewcommand{\headrulewidth}{0pt}
\renewcommand{\footrulewidth}{0pt}

\fancypagestyle{plain}{
    \fancyhf{}
    \fancyfoot[C]{\small NMIET, Department of Computer Engineering 2025-26}
    \fancyfoot[R]{\small \thepage}
    \renewcommand{\headrulewidth}{0pt}
    \renewcommand{\footrulewidth}{0pt}
}

% ─── Chapter Title Format ───────────────────────────────────
\titleformat{\chapter}[display]
  {\normalfont\Large\bfseries\centering}
 {CHAPTER \Roman{chapter}\vspace{0.15cm}\hrule}
  {5pt}
  {\Large\centering}
\titlespacing*{\chapter}{0pt}{-1cm}{8pt}

\titleformat{\section}{\normalfont\large\bfseries}{\thesection}{1em}{}
\titleformat{\subsection}{\normalfont\normalsize\bfseries}{\thesubsection}{1em}{}

% ════════════════════════════════════════════════════════════
\begin{document}

\AddToShipoutPictureBG{%
    \begin{tikzpicture}[overlay, remember picture]
        \draw[line width=1.5pt]
            ($ (current page.north west) + (1cm,-1cm) $)
            rectangle
            ($ (current page.south east) + (-1cm,1cm) $);
        \draw[line width=0.4pt]
            ($ (current page.north west) + (1.15cm,-1.15cm) $)
            rectangle
            ($ (current page.south east) + (-1.15cm,1.15cm) $);
    \end{tikzpicture}%
}

% ════════════════════════════════════════════════════════════
%  PAGE 1 — TITLE PAGE
% ════════════════════════════════════════════════════════════
\begin{titlepage}
\thispagestyle{fancy}
\pagenumbering{gobble}
\centering

\vspace*{-1cm}
{\large An}\\[0.3cm]
{\LARGE \textit{Internship Report}}\\[0.5cm]

{\normalsize SUBMITTED TO THE SAVITRIBAI PHULE PUNE UNIVERSITY, PUNE}\\
{\normalsize IN THE PARTIAL FULFILLMENT OF THE REQUIREMENTS}\\
{\normalsize FOR THE AWARD OF THE DEGREE}\\[0.4cm]

{\normalsize \textit{of}}\\[0.3cm]

{\large \textbf{BACHELOR OF ENGINEERING}}\\[0.3cm]
{\normalsize \textit{in}}\\[0.3cm]
{\large \textbf{COMPUTER ENGINEERING}}\\[0.8cm]

{\normalsize SUBMITTED BY}\\
{\normalsize \textbf{ADITYA SHAILESH BHARATE}}\\[0.8cm]

{\normalsize \textbf{UNDER THE GUIDANCE OF}}\\
{\normalsize PROF -- ASHWINI UTTURE}\\[0.8cm]

{\normalsize DEPARTMENT OF COMPUTER ENGINEERING}\\
{\normalsize (NBA ACCREDITED)}\\[0.8cm]

% ── Replace 'nmiet_logo.png' with your actual logo file ──
\includegraphics[width=0.22\textwidth]{nmiet_logo.png}\\[0.8cm]

{\normalsize NUTAN MAHARASHTRA INSTITUTE OF ENGINEERING \& TECHNOLOGY}\\
{\normalsize \textbf{TALEGAON DABHADE TAL.MAVAL, PUNE, 410507}}\\
{\normalsize SAVITRIBAI PHULE PUNE UNIVERSITY}\\
{\normalsize 2025 --2026}

\end{titlepage}


% ════════════════════════════════════════════════════════════
%  PAGE 2 — CERTIFICATE
% ════════════════════════════════════════════════════════════
\pagenumbering{gobble}
\chapter*{}
\thispagestyle{fancy}
\vspace*{-1.5cm}
\begin{center}
    {\normalsize PCET's \& NMVPM's}\\[0.2cm]
    {\large \textbf{NUTAN MAHARASHTRA INSTITUTE OF ENGINEERING \& TECHNOLOGY}}\\
    {\normalsize TALEGAON DABHADE, PUNE -- 410507}\\[0.5cm]

    % ── Replace 'nmiet_logo.png' with your actual logo file ──
    \includegraphics[width=0.22\textwidth]{nmiet_logo.png}\\[0.5cm]

    {\Huge \textit{Certificate}}\\[0.8cm]
\end{center}

\noindent This is to certify that, \textbf{Aditya Shailesh Bharate (Roll No. 01)} have successfully 
completed the Internship entitled and ``\underline{AWS AI-POWERED CLOUD ENGINEER}'' 
under my supervision, in the partial fulfilment of Bachelor of Engineering in Computer 
Engineering of Savitribai Phule Pune University (SPPU).

\vspace{2cm}
\noindent Date: 20/04/2026 \hfill Place: Pune

\vspace{3cm}
\noindent
\vspace{1.5cm}

\begin{tabular*}{\textwidth}{@{\extracolsep{\fill}}ccc}
    Prof. Ashwini Utture & Prof. Nilesh Kamble & Dr. Rahul Patil \\[0.2cm]
    \textbf{Internship Guide} & \textbf{Internship Coordinator} & \textbf{HOD} \\
    & & \textbf{Computer Engineering}
\end{tabular*}


% ════════════════════════════════════════════════════════════
%  PAGE 3 — ACKNOWLEDGEMENT
% ════════════════════════════════════════════════════════════
\chapter*{ACKNOWLEDGEMENT}
\addcontentsline{toc}{chapter}{Acknowledgement}
\thispagestyle{fancy}

I would like to express our deepest gratitude and appreciation to all those who have 
contributed to the successful completion of this Internship. It is with great pleasure that we 
acknowledge their invaluable support, guidance, and encouragement throughout this journey.

I thank our Internship guide \textbf{Prof. Ashwini Utture}, for her exceptional support, 
guidance, and patience throughout this journey. Her expertise and mentorship have been 
invaluable in shaping our understanding of the internship and overcoming challenges.

I would like to express our gratitude to our Internship coordinator \textbf{Prof. Nilesh Kamble} 
for their valuable suggestions, encouragement, and motivation that have contributed to the 
successful completion of this internship. We extend our appreciation to all the faculty members 
for their continuous support and encouragement during the course of this Internship.

I express our heartily gratitude towards \textbf{Dr. Rahul Patil}, Head of Department of 
Computer Engineering for guiding us to understand the work conceptually and also for 
providing necessary information and required resources with his constant encouragement to 
complete this internship work.

With deep sense of gratitude, I thank our principal \textbf{Dr. Pramod Patil} and Management 
of the NMIET for providing all necessary facilities and their constant encouragement and 
support. The encouragement and academic environment provided by our institution were 
crucial in fostering our creativity and enabling us to explore new ideas.

My sincere thanks to all the teaching and non-teaching staff members of the Computer 
Engineering Department for their assistance, cooperation, and provision of the required 
resources throughout this Internship. I am ending this acknowledgement with deep 
indebtedness to my friends who have helped us.

I sincerely acknowledge and appreciate the contributions of all those who directly or 
indirectly supported us in the completion of this Internship.


% ════════════════════════════════════════════════════════════
%  PAGE 4 — TABLE OF CONTENTS
% ════════════════════════════════════════════════════════════
\tableofcontents
\vspace*{-2cm}
\thispagestyle{fancy}


% ════════════════════════════════════════════════════════════
%  CHAPTER I — ORGANIZATIONAL STRUCTURE OF INDUSTRY
% ════════════════════════════════════════════════════════════
\clearpage
\pagenumbering{arabic}
\setcounter{page}{1}
\chapter{ORGANIZATIONAL STRUCTURE OF INDUSTRY}

\section{Organizational Structure:}
The organizational structure of Eduskill Foundation follows a hierarchical model with top management overseeing strategic decisions and departments such as development, data science, and operations handling execution. Team leads guide developers and interns in project tasks, ensuring efficient workflow and collaboration within the organization.

% If you have an org chart image, insert it here:
% \begin{figure}[h!]
%     \centering
%     \includegraphics[width=0.7\textwidth]{org_chart.png}
%     \caption{Organizational Structure of Eduskills Foundation}
% \end{figure}

\section{Address of the Company:}

\noindent \textbf{Corp. Office:} B-7, Awfis - Allied House, Near AICTE, Nelson Mandela Marg, 
Vasant Kunj, New Delhi -- 110070

\noindent \textbf{Head Office:} \#806, DLF Cyber City, Technology Corridor, 
Bhubaneswar - 751024, Odisha.

\noindent \textbf{Name:} Eduskills Foundation

\noindent \textbf{CIN:} [Enter CIN Number]

\section{Company Vision:}

Transforming the vision of skilled India and education for to benefit the education 
ecosystem by providing 360 degree holistic solutions to all the stakeholders.

\section{Company Mission:}

To positively impact 1 million beneficiaries by 2026.
By comprehensive identification of skills gaps in the students and mapping them with 
latest and world's best technical skills.


% ════════════════════════════════════════════════════════════
%  CHAPTER II — TECHNOLOGY LEARNED IN COMPANY
% ════════════════════════════════════════════════════════════
\chapter{TECHNOLOGY LEARNED IN COMPANY}

During the internship conducted through the Amazon Web Services Educate platform, 
several modern technologies related to cloud computing and artificial intelligence 
were studied. The training focused on understanding how cloud platforms help 
organizations store data, run applications, and manage infrastructure efficiently without 
relying on physical hardware. This approach improves scalability, cost efficiency, and 
system reliability, making cloud computing a key technology used by modern 
businesses and technology companies.

The internship also introduced practical knowledge of different AWS services used for 
computing, storage, networking, databases, serverless applications, and artificial 
intelligence. Through structured learning modules and assessments, a strong foundation 
was developed in cloud infrastructure, serverless architecture, machine learning 
fundamentals, and generative AI technologies, helping to understand how intelligent 
applications can be built using cloud platforms.

\medskip
\noindent \textbf{Technologies Learned}

\begin{itemize}[leftmargin=1.5cm]
    \item Cloud Computing Fundamentals
    \item Amazon S3 -- Cloud Storage Service
    \item Amazon EC2 -- Virtual Cloud Computing Instances
    \item Virtual Private Cloud (VPC) -- Cloud Networking
    \item Amazon RDS -- Managed Relational Database Service
    \item AWS Lambda -- Serverless Computing
    \item Amazon SageMaker -- Machine Learning Platform
    \item Amazon Bedrock -- Generative AI Service
    \item Cloud Security and Monitoring Concepts
    \item AI-Powered Cloud Application Development
\end{itemize}


% ════════════════════════════════════════════════════════════
%  CHAPTER III — DESCRIPTION OF INTERNSHIP
% ════════════════════════════════════════════════════════════
\chapter{DESCRIPTION OF INTERNSHIP}

During my internship, I gained foundational knowledge and practical exposure to cloud 
computing and artificial intelligence technologies through the Amazon Web Services 
Educate learning platform. The training consisted of structured modules designed to 
introduce key cloud concepts and widely used AWS services. Through these modules, I 
explored how cloud platforms enable organizations to store data, deploy applications, 
and manage infrastructure efficiently without relying on traditional on-premise systems.

The internship was organized into multiple learning modules covering topics such as 
cloud fundamentals, storage services, compute services, networking, databases, cloud 
security, serverless computing, machine learning, and generative AI. Each module 
included instructional materials, demonstrations, and assessments that helped me 
understand the working principles of different cloud services and their role in building 
modern cloud-based applications.

During the training, I explored several important AWS services such as Amazon S3 for 
cloud storage, Amazon EC2 for virtual computing resources, and AWS Lambda for 
serverless computing. Learning about these technologies helped me understand how 
developers can build applications that automatically scale based on demand without 
managing physical infrastructure. I also studied database management using Amazon 
RDS, which simplifies database deployment and maintenance in cloud environments.

Additionally, the internship introduced emerging technologies such as machine 
learning and generative artificial intelligence using tools like Amazon SageMaker 
and Amazon Bedrock. These modules provided insight into how AI models can be 
trained, deployed, and integrated with cloud services to develop intelligent applications. 
Overall, the internship broadened my understanding of modern cloud ecosystems and 
gave me practical exposure to technologies that are widely used in today's technology 
industry.


% ════════════════════════════════════════════════════════════
%  CHAPTER IV — DESCRIPTION OF TASKS
% ════════════════════════════════════════════════════════════
\chapter{DESCRIPTION OF TASKS}

During the internship, I completed a structured set of learning modules through the 
Amazon Web Services Educate platform as part of the AI Powered Cloud Engineer 
training program. Each module focused on a different aspect of cloud computing and 
artificial intelligence technologies that are widely used in modern IT infrastructure.

The initial module introduced the fundamentals of cloud computing, including the 
concept of cloud infrastructure, benefits of cloud adoption, and different service models 
used by cloud providers. The next modules focused on cloud storage and computing 
services, where I studied how Amazon S3 is used for storing and managing large 
amounts of data and how Amazon EC2 provides scalable virtual computing resources.

Further modules covered cloud networking and databases, where I explored the concept 
of Virtual Private Cloud (VPC) and learned how organizations manage relational 
databases using Amazon RDS. Additional modules introduced cloud operations and 
security, focusing on monitoring systems, access management, and protection of cloud 
resources.

The internship also included advanced topics such as serverless computing using AWS 
Lambda, machine learning foundations using Amazon SageMaker, and generative 
artificial intelligence using Amazon Bedrock. These modules helped me understand how 
cloud platforms can support intelligent and scalable applications.
\newpage
\section{Introduction of Projects:}

\noindent\textbf{\large Event-Driven Serverless System for AWS Cost Optimization}

\vspace{0.3cm}
\noindent\textbf{Introduction}

This project focuses on optimizing cloud costs by automatically identifying and stopping 
idle Amazon EC2 instances. In cloud environments, instances often remain running even 
when not actively used, resulting in unnecessary expenses. Manual monitoring is 
inefficient and difficult to scale, creating a need for an automated solution.

\noindent\textbf{Proposed Solution}

An event-driven serverless architecture was implemented using AWS services. Amazon 
EventBridge triggers an AWS Lambda function at scheduled intervals. The Lambda 
function retrieves EC2 instance details and evaluates their CPU utilization using Amazon 
CloudWatch metrics. Instances with CPU usage below a defined threshold are considered 
idle and are automatically stopped to reduce costs.

\noindent\textbf{Key Features}
\begin{itemize}[leftmargin=1.5cm]
    \item Automated detection and stopping of idle EC2 instances
    \item Use of CloudWatch metrics for performance evaluation
    \item Tag-based filtering to target only non-production resources
    \item Logging mechanism for monitoring and tracking system actions
    \item Fully serverless and scalable architecture
\end{itemize}

\noindent\textbf{Challenges and Approach}

A key challenge was defining what qualifies as an ``idle'' instance. This was addressed 
by setting appropriate CPU utilization thresholds based on usage patterns, ensuring 
accurate identification without affecting necessary workloads.

\noindent\textbf{Conclusion}

The project provides a scalable and efficient solution for cloud cost optimization. By 
leveraging AWS serverless services, it reduces manual effort, improves resource 
utilization, and minimizes unnecessary cloud spending while maintaining operational 
reliability.

\newpage
\section{Motivation and Relational Study:}

The increasing use of digital platforms and online services has created a strong demand 
for scalable computing infrastructure. Cloud computing has emerged as a powerful 
solution that allows organizations to access computing resources such as storage, 
processing power, and databases through the internet without maintaining their own 
hardware systems.

The motivation behind undertaking this internship was to gain practical exposure to 
cloud technologies and understand how modern IT systems are designed using cloud 
platforms. Learning about cloud computing provides valuable insight into how 
businesses manage data, deploy applications, and handle large-scale workloads 
efficiently.

Another important motivation was to explore the integration of artificial intelligence 
with cloud computing. AI technologies such as machine learning and generative AI 
require large amounts of data and computing resources, which are efficiently supported 
by cloud platforms. Understanding how these technologies interact helps in developing 
modern intelligent systems.

This internship also helped bridge the gap between theoretical knowledge and real-world 
technology applications. Concepts learned in areas such as networking, databases, and 
software systems could be related to practical cloud services provided by AWS. By 
studying these technologies together, I developed a clearer understanding of how cloud 
infrastructure, data processing systems, and AI models work collectively to support 
modern digital applications.
\newpage
\section{Methodology Details:}

The methodology followed during the internship was based on a structured learning 
approach through the Amazon Web Services Educate training modules. Each module 
focused on a specific cloud technology and included learning materials and assessments 
to evaluate understanding. This step-by-step approach helped build a clear 
understanding of how different cloud services operate and how they are used in 
real-world applications.

The process began with studying the fundamentals of cloud computing, including cloud 
architecture, service models, and the benefits of cloud adoption such as scalability, 
reliability, and cost efficiency. I also learned about different cloud service models like 
Infrastructure as a Service (IaaS), Platform as a Service (PaaS), and Software as a 
Service (SaaS).

The next modules focused on cloud storage using Amazon S3, where I learned how 
large volumes of data can be securely stored and managed in the cloud. This was 
followed by studying cloud computing resources through Amazon EC2, which 
explained how virtual machines are deployed and scaled based on application 
requirements.

Further modules introduced cloud networking through Virtual Private Cloud (VPC) 
and database management using Amazon RDS, which demonstrated how secure 
networks and relational databases are managed in cloud environments.

Another important module introduced serverless computing using AWS Lambda. 
Serverless computing allows developers to run application code without managing 
underlying servers. The module explained how event-driven architecture works and how 
applications can automatically scale based on incoming requests. This concept 
highlighted the efficiency and flexibility of modern cloud-based application development.

The final modules focused on machine learning and artificial intelligence technologies. 
Using Amazon SageMaker, I studied the basic workflow of machine learning including 
data preparation, model training, and deployment. Additionally, the internship 
introduced generative artificial intelligence using Amazon Bedrock, which provides 
access to advanced AI models capable of generating text and other forms of content.


% ════════════════════════════════════════════════════════════
%  CHAPTER V — LEARNING OUTCOMES
% ════════════════════════════════════════════════════════════
\chapter{LEARNING OUTCOMES}

Through this internship, I gained a comprehensive understanding of the fundamental concepts of cloud computing and artificial intelligence technologies. The training provided deep insights into how cloud platforms deliver scalable and on-demand infrastructure for storing data, running applications, and efficiently managing computing resources. I also developed a clear understanding of various cloud service models such as Infrastructure as a Service (IaaS), Platform as a Service (PaaS), and Software as a Service (SaaS), and how organizations leverage these models to enhance performance, reliability, and cost efficiency.

During the internship, I acquired practical knowledge of several core services offered by Amazon Web Services (AWS). I learned how Amazon S3 is utilized for secure, durable, and scalable object storage, enabling efficient data management. Additionally, I gained hands-on experience with Amazon EC2, which provides virtual servers for deploying and managing applications in a flexible cloud environment. I also explored cloud networking concepts, including Virtual Private Cloud (VPC), subnets, and security groups, which are essential for building secure and well-structured cloud architectures.

Furthermore, I developed an understanding of database management using Amazon RDS, which simplifies the setup, operation, and scaling of relational databases in the cloud. I also gained exposure to serverless computing through AWS Lambda, where I learned how to execute code without managing servers, enabling cost-effective and event-driven application development.


Moreover, this internship enhanced my problem-solving abilities and technical skills by allowing me to work on real-world use cases, such as optimizing cloud costs through automation. It also improved my understanding of best practices in cloud architecture, resource management, and system design.

Overall, the internship played a crucial role in strengthening my technical foundation and provided valuable exposure to modern cloud and AI technologies. It has equipped me with the knowledge and skills required to design scalable, efficient, and intelligent applications by integrating cloud computing with artificial intelligence.


% ════════════════════════════════════════════════════════════
%  CHAPTER VI — CONCLUSION
% ════════════════════════════════════════════════════════════
\chapter{CONCLUSION}

The internship provided valuable exposure to the fundamentals of cloud computing and 
artificial intelligence technologies through the learning modules on the Amazon Web 
Services Educate platform. Throughout the internship, I studied various cloud services 
and gained an understanding of how modern cloud infrastructure supports the 
development and deployment of scalable applications. The training modules helped me 
understand the architecture and working principles of different cloud technologies.

During the internship, I explored important AWS services such as Amazon S3 for cloud 
storage, Amazon EC2 for computing resources, Amazon RDS for database management, 
and AWS Lambda for serverless computing. I also gained introductory knowledge of 
machine learning and generative AI using Amazon SageMaker and Amazon Bedrock. 
These technologies demonstrated how cloud platforms enable efficient data processing, 
application deployment, and intelligent system development.

Overall, the internship strengthened my understanding of cloud-based technologies and 
provided insight into how cloud computing and artificial intelligence are used in modern 
IT environments. The knowledge gained during this training will help in further 
developing technical skills and exploring advanced technologies in the field of cloud 
computing and AI.


% ════════════════════════════════════════════════════════════
%  CHAPTER VII — REFERENCES
% ════════════════════════════════════════════════════════════
\chapter*{\vspace{-1cm}REFERENCES}
\addcontentsline{toc}{chapter}{References}
\thispagestyle{fancy}

\begin{center}
    \rule{\linewidth}{0.4pt}
\end{center}

\vspace{0.4cm}

\begin{enumerate}[leftmargin=1.5cm, label={[\arabic*]}, itemsep=0.3cm]
\justifying

    \item Amazon Web Services, \textit{AWS Lambda -- Serverless Computing Documentation}. [Online]. Available: \url{https://docs.aws.amazon.com/lambda/}

    \item Amazon Web Services, \textit{Amazon SageMaker -- Machine Learning Documentation}. [Online]. Available: \\ \url{https://docs.aws.amazon.com/sagemaker/}

    \item Amazon Web Services, \textit{Amazon Bedrock -- Generative AI Documentation}. [Online]. Available: \url{https://docs.aws.amazon.com/bedrock/}

    \item Amazon Web Services, \textit{Amazon EventBridge Documentation}. [Online]. Available: \url{https://docs.aws.amazon.com/eventbridge/}

    \item Amazon Web Services, \textit{Amazon CloudWatch -- Monitoring and Observability}. [Online]. Available: \\ \url{https://docs.aws.amazon.com/cloudwatch/}

    \item E. Jonas et al., ``Cloud programming simplified: A Berkeley view on serverless computing,'' \textit{arXiv preprint}, arXiv:1902.03383, Feb. 2019. [Online]. Available: \url{https://arxiv.org/abs/1902.03383}

    \item P. P. Singh and K. S. Saini, ``Optimizing AWS Cloud Infrastructure: An Analysis of Cloud-Nuke's Impact on Productivity, Cost Efficiency, and Resource Management,'' in \textit{Proc. 2024 3rd International Conference on Sentiment Analysis and Deep Learning (ICSADL)}, 2024, doi: 10.1109/ICSADL62624.2024.10601474.

   

    \item K. Ren, C. Wang, and Q. Wang, ``Security challenges for the public cloud,'' \textit{IEEE Internet Computing}, vol. 16, no. 1, pp. 69--73, Jan.--Feb. 2012, doi: 10.1109/MIC.2012.14.

\end{enumerate}

% ════════════════════════════════════════════════════════════
%  ANNEXURE B — COMPLETION CERTIFICATE
% ════════════════════════════════════════════════════════════
\appendix
\chapter*{\vspace{-1cm}Annexure B}
\addcontentsline{toc}{chapter}{Annexure B -- Completion Certificate}
\thispagestyle{fancy}

\vspace{0.5cm}

\begin{center}
    {\Large \textbf{COMPLETION CERTIFICATE}}
\end{center}

\vspace{0.3cm}

\begin{center}
\includegraphics[width=0.95\textwidth, height=0.75\textheight, keepaspectratio]{completion_certificate.png}
\end{center}

\vspace{0.3cm}


% ════════════════════════════════════════════════════════════
%  LOG BOOK
% ════════════════════════════════════════════════════════════
\chapter*{Log Book}
\addcontentsline{toc}{chapter}{Log Book}
\thispagestyle{fancy}

\begin{center}
\small
\begin{tabular}{|p{1.2cm}|p{9cm}|p{2.8cm}|}
    \hline
    \textbf{Week} & \textbf{Work Done} & \textbf{Coordinator Sign} \\
    \hline
    Week 1 & Studied Cloud Computing Fundamentals -- cloud architecture, service models (IaaS, PaaS, SaaS) and benefits of cloud adoption via AWS Educate. & \\[0.4cm]
    \hline
    Week 2 & Explored Amazon S3 -- cloud storage concepts, bucket creation, object storage and access policies. & \\[0.4cm]
    \hline
    Week 3 & Studied Amazon EC2 -- virtual computing instances, instance types, launch configurations and application deployment. & \\[0.4cm]
    \hline
    Week 4 & Learned VPC and Amazon RDS -- cloud networking, subnets, security groups and managed relational database deployment. & \\[0.4cm]
    \hline
    Week 5 & Studied Cloud Security and Monitoring -- IAM roles, access management, CloudWatch metrics, logging and protection of cloud resources. & \\[0.4cm]
    \hline
    Week 6 & Explored AWS Lambda and Serverless Computing -- event-driven architecture, function deployment and automatic scaling concepts. & \\[0.4cm]
    \hline
    Week 7 & Studied Amazon SageMaker -- machine learning workflow including data preparation, model training, evaluation and cloud deployment. & \\[0.4cm]
    \hline
    Week 8 & Explored Amazon Bedrock and Generative AI -- large language models, generative AI capabilities and completed all module assessments. & \\[0.4cm]
    \hline
    Week 9 & Built an event-driven serverless system on AWS using EventBridge and Lambda to auto-stop idle EC2 instances based on CloudWatch CPU metrics. Implemented tag-based filtering for non-production resources and CPU threshold-based idle detection to reduce cloud costs. & \\[0.4cm]
    \hline
\end{tabular}
\end{center}

\end{document}
